\section*{Chapitre \ref{ch:b3:electronic-spectroscopy} --- Spectroscopie électronique et photophysique}

\begin{solution}{exo:b3:electronic-spectroscopy:1}
(a) Interdite par la règle de spin ($\Delta S = 1$). (b) Interdite par la
\omterm{thm:b3:electronic-spectroscopy:laporte}{règle de Laporte}
($g \to g$)~; observée faiblement grâce au couplage vibronique.
(c) Permise~: $B_{1u}$ est la
représentation de $z$ dans $D_{2h}$. (d) Interdite par la symétrie~:
$A_2$ n'est la symétrie d'aucune des coordonnées $x$,
$y$, $z$ dans $C_{2v}$~; la bande est faible.
\end{solution}

\begin{solution}{exo:b3:electronic-spectroscopy:2}
$10^7/349 - 10^7/450 = 28653 - 22222 = \qty{6431}{cm^{-1}}$, \qty{0.80}{eV}.
\end{solution}

\begin{solution}{exo:b3:electronic-spectroscopy:3}
$A = 5700 \times 1 \times 1.0\times10^{-5} = 0.057$~; fraction absorbée $1 - 10^{-0.057} =
0.12$.
\end{solution}

\begin{solution}{exo:b3:electronic-spectroscopy:4}
$k_r = \Phi_F/\tau = \qty{1.5e8}{s^{-1}}$, $\tau_r = \qty{6.7}{ns}$~; $k_{nr} = (1 - \Phi_F)/\tau
= \qty{1.0e8}{s^{-1}}$.
\end{solution}

\begin{solution}{exo:b3:electronic-spectroscopy:5}
$\int\varepsilon\,\dd\tilde\nu = 1.0645 \times 5700 \times 5000 = \num{3.03e7}$, donc $f \approx
4.32\times10^{-9} \times 3.03\times10^7 = 0.13$~: une transition permise
d'intensité modérée.
\end{solution}

\begin{solution}{exo:b3:electronic-spectroscopy:6}
Le facteur de Poisson $\eu^{-S}S^v/v!$ est maximal pour $v = \lfloor S\rfloor$. $S = 0.3$~: la
raie 0--0 elle-même. $S = 1.5$~: $v' = 1$, 1.5 fois la raie 0--0. $S = 5$~:
$v' = 4$
et 5 à égalité, 26 fois la raie 0--0.
\end{solution}

\begin{solution}{exo:b3:electronic-spectroscopy:7}
$\sum k = \qty{4.2e8}{s^{-1}}$~: $\tau = \qty{2.4}{ns}$, $\Phi_F = 0.24$, $\Phi_{\mathrm{ISC}} = 0.71$
(et $\Phi_{\mathrm{IC}} = 0.05$).
\end{solution}

\begin{solution}{exo:b3:electronic-spectroscopy:8}
$\Phi = 0.546 \times 0.46 \times (0.050/0.040) \times (1.4266/1.333)^2 = 0.36$.
\end{solution}

\begin{solution}{exo:b3:electronic-spectroscopy:9}
$\tau_0/\tau = 1.429$ et 1.860~: linéaire en $[Q]$, donc c'est la durée de
vie elle-même qui est raccourcie~: \omterm{def:b3:electronic-spectroscopy:quencher}{inhibition dynamique}. $K_{SV} = 43$
L/mol (les deux points), $k_q = 43/8.0
\times 10^{-9} = \qty{5.4e9}{L.mol^{-1}.s^{-1}}$.
\end{solution}

\begin{solution}{exo:b3:electronic-spectroscopy:10}
\omterm{def:b3:electronic-spectroscopy:quencher}{Inhibition statique}~: les molécules qui émettent encore vivent aussi
longtemps qu'avant~; la moitié des molécules sont liées, dans l'état
fondamental, en un complexe non fluorescent. Avec une constante
d'association $K_a$, la fraction libre vaut $1/(1 + K_a[Q])$ et $I_0/I = 1 +
K_a[Q]$, la même forme que l'\omterm{thm:b3:electronic-spectroscopy:stern-volmer}{équation de Stern--Volmer}, mais avec $\tau$
inchangée.
\end{solution}

\begin{solution}{exo:b3:electronic-spectroscopy:11}
$\tau_r = \tau/\Phi_F = \qty{14}{ns}$~; $\Phi_T = 1 - 0.36 = 0.64$. La
durée de vie mesurée est raccourcie par le \omterm{def:b3:electronic-spectroscopy:radiationless}{croisement intersystème}
concurrent~; $\tau_r$ est ce que donnerait l'émission seule.
\end{solution}

\begin{solution}{exo:b3:electronic-spectroscopy:12}
$\frac12\langle\alpha\beta - \beta\alpha|\alpha\beta + \beta\alpha\rangle = \frac12(1 + 0 - 0 - 1) = 0$, en utilisant
$\langle\alpha|\alpha\rangle = \langle\beta|\beta\rangle = 1$, $\langle\alpha|\beta\rangle = 0$ pour chaque électron.
L'\omterm{def:b3:quantum-model-systems:operator}{opérateur} dipolaire n'agit pas sur le spin, donc le moment de
transition contient ce facteur nul~: $T_1 \leftarrow S_0$ n'a aucune
intensité, sauf si le \omterm{def:b3:many-electron-atoms:spin-orbit}{couplage spin--orbite} mélange
les états.
\end{solution}

\begin{solution}{pb:b3:electronic-spectroscopy:1}
\textbf{1.} $A = 5700 \times 1 \times 1.0\times10^{-4} = 0.57$.
\textbf{2.} $1 - 10^{-0.57} = 0.73$.
\textbf{3.} Un $\varepsilon$ de quelques milliers~: une transition
$\pi\to\pi^*$ permise, d'intensité modérée.
\textbf{4.} $1239.8/349 = \qty{3.55}{eV}$.
\textbf{5.} Elle n'absorbe que dans l'ultraviolet~; la lumière visible
passe.
\textbf{6.} Pour que la lumière absorbée reste proportionnelle à la
concentration et pour éviter la réabsorption de la lumière émise (effets
de filtre interne).
\textbf{7.} $S_0 \to S_1$ (absorption, puis \omterm{def:b3:electronic-spectroscopy:radiationless}{relaxation vibrationnelle})~;
depuis $S_1$~: \omterm{def:b3:electronic-spectroscopy:luminescence}{fluorescence}, \omterm{def:b3:electronic-spectroscopy:radiationless}{conversion interne}, \omterm{def:b3:electronic-spectroscopy:radiationless}{croisement intersystème}
vers $T_1$.
\textbf{8.} $k_r = 0.546/19 \times 10^{-9} = \qty{2.87e7}{s^{-1}}$~; $\tau_r = \qty{35}{ns}$.
\textbf{9.} $(1 - 0.546)/\tau_0 = \qty{2.39e7}{s^{-1}}$.
\textbf{10.} $28653 - 22222 = \qty{6431}{cm^{-1}}$.
\textbf{11.} L'émission part de $S_1$ relaxé et aboutit sur des niveaux
de $S_0$ excités en vibration, après relaxation de la molécule et du
solvant~: le photon a moins d'énergie, ici assez peu pour tomber dans le
visible.
\textbf{12.} $\Phi_F \times (349/450) = 0.42$~: 42\,\% de l'énergie
absorbée.
\textbf{13.} Les points croissent de 0.594 par \qty{0.010}{mol/L}~: une
droite passant par 1.
\textbf{14.} Moindres carrés sur les cinq points~: pente $K_{SV} = \qty{59.4}{L/mol}$,
ordonnée à l'origine 1.000.
\textbf{15.} $k_q = K_{SV}/\tau_0 = \qty{3.1e9}{L.mol^{-1}.s^{-1}}$.
\textbf{16.} $\tau = \tau_0/(1 + 59.4 \times 0.040) = \qty{5.6}{ns}$.
\textbf{17.} Mesurer les durées de vie~: pour une \omterm{def:b3:electronic-spectroscopy:quencher}{inhibition dynamique},
$\tau_0/\tau$ se place sur la même droite que $I_0/I$.
\textbf{18.} Le tonic ne contient pas de chlorure pour l'inhiber, et il
est assez acide pour maintenir la quinine protonée, sa forme
fluorescente.
\textbf{19.} $k_d = 8 \times 8.314 \times 298/(3 \times 0.890\times10^{-3}) =
\qty{7.4e6}{m^3.mol^{-1}.s^{-1}} = \qty{7.4e9}{L.mol^{-1}.s^{-1}}$.
\textbf{20.} $k_q$ vaut un peu moins de la moitié de $k_d$.
\textbf{21.} Environ 42\,\%.
\textbf{22.} Plus lente~: $k_d \propto 1/\eta$, et $k_q$ ne peut que
diminuer avec elle.
\textbf{23.} \textbf{$k_q/k_d = 0.42$}~: le chlorure inhibe la quinine à
près de la moitié de la vitesse des rencontres contrôlées par la
diffusion.
\end{solution}
